\documentclass[10pt,a4paper]{amsart}
\usepackage{amsmath,amssymb,amsthm,graphicx}
% Retained minimal preamble for the Stacks Project chapter varieties.tex.
% Verbatim excerpt of the chapter preamble: Xy-pic setup (native preamble lines 14-18), the native theorem
% declarations with the mathematical macros (native preamble lines 154-217), and the project own no-op
% definitions of the reference/slogan/history display environments (native wrapper.tex lines 1-3) that the
% counted statement uses for its slogan. Copied mechanically; nothing added, removed or redefined.
% For commutative diagrams we use Xy-pic
\usepackage[all]{xy}

% We use 2cell for 2-commutative diagrams.
\xyoption{2cell}
% Theorem environments.
%
\theoremstyle{plain}
\newtheorem{theorem}[subsection]{Theorem}
\newtheorem{proposition}[subsection]{Proposition}
\newtheorem{lemma}[subsection]{Lemma}

\theoremstyle{definition}
\newtheorem{definition}[subsection]{Definition}
\newtheorem{example}[subsection]{Example}
\newtheorem{exercise}[subsection]{Exercise}
\newtheorem{situation}[subsection]{Situation}

\theoremstyle{remark}
\newtheorem{remark}[subsection]{Remark}
\newtheorem{remarks}[subsection]{Remarks}

\numberwithin{equation}{subsection}

% Macros
%
\def\lim{\mathop{\mathrm{lim}}\nolimits}
\def\colim{\mathop{\mathrm{colim}}\nolimits}
\def\Spec{\mathop{\mathrm{Spec}}}
\def\Hom{\mathop{\mathrm{Hom}}\nolimits}
\def\Ext{\mathop{\mathrm{Ext}}\nolimits}
\def\SheafHom{\mathop{\mathcal{H}\!\mathit{om}}\nolimits}
\def\SheafExt{\mathop{\mathcal{E}\!\mathit{xt}}\nolimits}
\def\Sch{\mathit{Sch}}
\def\Mor{\mathop{\mathrm{Mor}}\nolimits}
\def\Ob{\mathop{\mathrm{Ob}}\nolimits}
\def\Sh{\mathop{\mathit{Sh}}\nolimits}
\def\NL{\mathop{N\!L}\nolimits}
\def\CH{\mathop{\mathrm{CH}}\nolimits}
\def\proetale{{pro\text{-}\acute{e}tale}}
\def\etale{{\acute{e}tale}}
\def\QCoh{\mathit{QCoh}}
\def\Ker{\mathop{\mathrm{Ker}}}
\def\Im{\mathop{\mathrm{Im}}}
\def\Coker{\mathop{\mathrm{Coker}}}
\def\Coim{\mathop{\mathrm{Coim}}}

% Boxtimes
%
\DeclareMathSymbol{\boxtimes}{\mathbin}{AMSa}{"02}

%
% Macros for moduli stacks/spaces
%
\def\QCohstack{\mathcal{QC}\!\mathit{oh}}
\def\Cohstack{\mathcal{C}\!\mathit{oh}}
\def\Spacesstack{\mathcal{S}\!\mathit{paces}}
\def\Quotfunctor{\mathrm{Quot}}
\def\Hilbfunctor{\mathrm{Hilb}}
\def\Curvesstack{\mathcal{C}\!\mathit{urves}}
\def\Polarizedstack{\mathcal{P}\!\mathit{olarized}}
\def\Complexesstack{\mathcal{C}\!\mathit{omplexes}}
% \Pic is the operator that assigns to X its picard group, usage \Pic(X)
% \Picardstack_{X/B} denotes the Picard stack of X over B
% \Picardfunctor_{X/B} denotes the Picard functor of X over B
\def\Pic{\mathop{\mathrm{Pic}}\nolimits}
\def\Picardstack{\mathcal{P}\!\mathit{ic}}
\def\Picardfunctor{\mathrm{Pic}}
\def\Deformationcategory{\mathcal{D}\!\mathit{ef}}
\newenvironment{reference}{}{}
\newenvironment{slogan}{}{}
\newenvironment{history}{}{}
\usepackage{cleveref}
\crefname{originallemma}{original lemma}{original lemmas}
\Crefname{originallemma}{Original lemma}{Original lemmas}
\makeatletter
\newlabel{lemma-geometrically-reduced-at-point}{{035W}{source}{}{sourcepointer-1}{}}
\newlabel{lemma-geometrically-reduced-at-point@cref}{{[originallemma][0][]035W}{[1][0][]source}{}{}{}}
\makeatother
\crefname{originallemma}{original lemma}{original lemmas}
\Crefname{originallemma}{Original lemma}{Original lemmas}
\makeatletter
\newlabel{lemma-geometrically-regular-at-point}{{038U}{source}{}{sourcepointer-2}{}}
\newlabel{lemma-geometrically-regular-at-point@cref}{{[originallemma][0][]038U}{[1][0][]source}{}{}{}}
\makeatother
\crefname{originallemma}{original lemma}{original lemmas}
\Crefname{originallemma}{Original lemma}{Original lemmas}
\makeatletter
\newlabel{lemma-normalization-locally-algebraic}{{0BXR}{source}{}{sourcepointer-3}{}}
\newlabel{lemma-normalization-locally-algebraic@cref}{{[originallemma][0][]0BXR}{[1][0][]source}{}{}{}}
\makeatother
\crefname{originallemma}{original lemma}{original lemmas}
\Crefname{originallemma}{Original lemma}{Original lemmas}
\makeatletter
\newlabel{more-algebra-lemma-regular-implies-fs}{{07EJ}{source}{}{sourcepointer-4}{}}
\newlabel{more-algebra-lemma-regular-implies-fs@cref}{{[originallemma][0][]07EJ}{[1][0][]source}{}{}{}}
\makeatother
\crefname{originaltheorem}{original theorem}{original theorems}
\Crefname{originaltheorem}{Original theorem}{Original theorems}
\makeatletter
\newlabel{more-algebra-theorem-regular-fs}{{07EL}{source}{}{sourcepointer-5}{}}
\newlabel{more-algebra-theorem-regular-fs@cref}{{[originaltheorem][0][]07EL}{[1][0][]source}{}{}{}}
\makeatother
\crefname{originallemma}{original lemma}{original lemmas}
\Crefname{originallemma}{Original lemma}{Original lemmas}
\makeatletter
\newlabel{properties-lemma-criterion-normal}{{0345}{source}{}{sourcepointer-6}{}}
\newlabel{properties-lemma-criterion-normal@cref}{{[originallemma][0][]0345}{[1][0][]source}{}{}{}}
\makeatother
\crefname{originalsection}{original section}{original sections}
\Crefname{originalsection}{Original section}{Original sections}
\makeatletter
\newlabel{section-algebraic-schemes}{{06LF}{source}{}{sourcepointer-7}{}}
\newlabel{section-algebraic-schemes@cref}{{[originalsection][0][]06LF}{[1][0][]source}{}{}{}}
\makeatother
\title{Geometric reducedness and geometric regularity of the normalization at a one dimensional point with separable residue fields}
\author{The Stacks Project Authors}
\date{stacks-project:varieties.tex:local-snapshot-20261011; GNU FDL 1.2 or later}
\begin{document}
\maketitle
\noindent\textit{Editorial note.} All delivered mathematical text is the author's own, copied from the pinned Stacks Project chapter file varieties.tex; no statement or proof step was rewritten, reordered or removed.\par\medskip
\noindent\textit{Editorial note.} The proof opens with the author's own 'We will use the results of Lemma lemma-normalization-locally-algebraic without further mention'; that sentence is reproduced verbatim and the results in question are navigation stubs, while every step that uses them is written out inside the counted span.\par\medskip
\noindent\textit{Editorial note.} The counted spans carry 7 navigation stubs and no bibliography citation.\par\medskip
\noindent\textit{Editorial note.} The preamble is the chapter's retained minimal preamble (sources/stacks-varieties/retained-preamble.tex lines 1-77), a verbatim excerpt of the native preamble.\par\medskip
\noindent\textit{Editorial note.} The label lemma-geometrically-normal-in-codim-1 is defined only in varieties.tex, so no cross-chapter restatement had to be weighed for this card.\par\medskip
\noindent\textit{Allowed standard prerequisites (editorial).}\begin{itemize}
\item Normalization of a locally algebraic scheme: finiteness and dimension behaviour. Same chapter, Lemma lemma-normalization-locally-algebraic, invoked at proof lines 5113-5114. Supplies the finiteness of the normalization and the one dimensionality of the local rings above x.
\item A normal one dimensional local ring is a discrete valuation ring. Properties, Lemma properties-lemma-criterion-normal, invoked at proof lines 5116-5118. Supplies the discrete valuation ring structure of the local ring of the normalization.
\item Regular local algebras with separable residue field are formally smooth, and formal smoothness equals geometric regularity. More on Algebra, Lemma more-algebra-lemma-regular-implies-fs and Theorem more-algebra-theorem-regular-fs, invoked at proof lines 5119-5126. Supplies geometric regularity of the normalization at the point above x.
\item Pointwise criteria for geometric reducedness and geometric regularity. Same chapter, Lemmas lemma-geometrically-reduced-at-point and lemma-geometrically-regular-at-point, invoked at proof lines 5127-5128 and 5135-5136. Supplies the passage between local rings and the geometric properties of the scheme at a point.
\end{itemize}\medskip
\ifdefined\hypertarget\hypertarget{sourcepointer-1}{}\fi
\noindent\textit{Source reference.} Original lemma 035W: sources/stacks-varieties/pinned-original-source.tex, source lines 354--354. Pointer outside the counted span, invoked once for context; the load-bearing argument of the counted statement is the author's proof reproduced in full in the proof span.\par\medskip
\ifdefined\hypertarget\hypertarget{sourcepointer-2}{}\fi
\noindent\textit{Source reference.} Original lemma 038U: sources/stacks-varieties/pinned-original-source.tex, source lines 2382--2382. Pointer outside the counted span, invoked once for context; the load-bearing argument of the counted statement is the author's proof reproduced in full in the proof span.\par\medskip
\ifdefined\hypertarget\hypertarget{sourcepointer-3}{}\fi
\noindent\textit{Source reference.} Original lemma 0BXR: sources/stacks-varieties/pinned-original-source.tex, source lines 4903--4903. Pointer outside the counted span, invoked once for context; the load-bearing argument of the counted statement is the author's proof reproduced in full in the proof span.\par\medskip
\ifdefined\hypertarget\hypertarget{sourcepointer-4}{}\fi
\noindent\textit{Source reference.} Original lemma 07EJ: sources/stacks-chow/tags.txt, source lines 9449--9449. Pointer outside the counted span, invoked once for context; the load-bearing argument of the counted statement is the author's proof reproduced in full in the proof span.\par\medskip
\ifdefined\hypertarget\hypertarget{sourcepointer-5}{}\fi
\noindent\textit{Source reference.} Original theorem 07EL: sources/stacks-chow/tags.txt, source lines 9451--9451. Pointer outside the counted span, invoked once for context; the load-bearing argument of the counted statement is the author's proof reproduced in full in the proof span.\par\medskip
\ifdefined\hypertarget\hypertarget{sourcepointer-6}{}\fi
\noindent\textit{Source reference.} Original lemma 0345: sources/stacks-chow/tags.txt, source lines 4080--4080. Pointer outside the counted span, invoked once for context; the load-bearing argument of the counted statement is the author's proof reproduced in full in the proof span.\par\medskip
\ifdefined\hypertarget\hypertarget{sourcepointer-7}{}\fi
\noindent\textit{Source reference.} Original section 06LF: sources/stacks-varieties/pinned-original-source.tex, source lines 3713--3713. Pointer outside the counted span, invoked once for context; the load-bearing argument of the counted statement is the author's proof reproduced in full in the proof span.\par\medskip
\begin{lemma}
\label{lemma-geometrically-normal-in-codim-1}
Let $k$ be a field. Let $X$ be a locally algebraic $k$-scheme.
Let $\nu : X^\nu \to X$ be the normalization of $X$.
Let $x \in X$ be a point such that (a) $\mathcal{O}_{X, x}$
is reduced, (b) $\dim(\mathcal{O}_{X, x}) = 1$, and (c)
for every $x' \in X^\nu$ with $\nu(x') = x$ the extension
$\kappa(x')/k$ is separable. Then $X$ is geometrically reduced at $x$
and $X^\nu$ is geometrically regular at $x'$ with $\nu(x') = x$.
\end{lemma}

\begin{proof}
We will use the results of Lemma \ref{lemma-normalization-locally-algebraic}
without further mention. Let $x' \in X^\nu$ be a point over $x$.
By dimension theory (Section \ref{section-algebraic-schemes}) we have
$\dim(\mathcal{O}_{X^\nu, x'}) = 1$. Since $X^\nu$ is normal, we
see that $\mathcal{O}_{X^\nu, x'}$ is a discrete valuation ring
(Properties, Lemma \ref{properties-lemma-criterion-normal}).
Thus $\mathcal{O}_{X^\nu, x'}$ is a regular local $k$-algebra
whose residue field is separable over $k$. Hence
$k \to \mathcal{O}_{X^\nu, x'}$ is formally smooth
in the $\mathfrak m_{x'}$-adic topology, see
More on Algebra, Lemma \ref{more-algebra-lemma-regular-implies-fs}.
Then $\mathcal{O}_{X^\nu, x'}$ is geometrically regular
over $k$ by
More on Algebra, Theorem \ref{more-algebra-theorem-regular-fs}.
Thus $X^\nu$ is geometrically regular at $x'$ by
Lemma \ref{lemma-geometrically-regular-at-point}.

\medskip\noindent
Since $\mathcal{O}_{X, x}$ is reduced, the family of maps
$\mathcal{O}_{X, x} \to \mathcal{O}_{X^\nu, x'}$ is injective.
Since $\mathcal{O}_{X^\nu, x'}$ is a geometrically reduced
$k$-algebra, it follows immediately that $\mathcal{O}_{X, x}$
is a geometrically reduced $k$-algebra. Hence $X$ is geometrically
reduced at $x$ by Lemma \ref{lemma-geometrically-reduced-at-point}.
\end{proof}

\end{document}
